\documentclass[conference]{IEEEtran}
\IEEEoverridecommandlockouts
% The preceding line is only needed to identify funding in the first footnote. If that is unneeded, please comment it out.
\usepackage{cite}
\usepackage{amsmath,amssymb,amsfonts}
\usepackage{algorithmic}
\usepackage{graphicx}
\usepackage{textcomp}
\usepackage{xcolor}
\usepackage{float}
\usepackage{subcaption}
\usepackage{algorithm}
\usepackage{algpseudocode}


% ====== 新增以下這段來支援中文 ======
\usepackage[AutoFakeBold=true, AutoFakeSlant=true]{xeCJK}\setmainfont{Times New Roman}
% 設定主要的繁體中文字體，您可以根據您的作業系統或編譯環境更改
%\setCJKmainfont{Noto Sans TC} % Windows 常用的字體
\setCJKmainfont{標楷體}   % 如果想要比較正式的襯線字體

\def\BibTeX{{\rm B\kern-.05em{\sc i\kern-.025em b}\kern-.08em
    T\kern-.1667em\lower.7ex\hbox{E}\kern-.125emX}}
\begin{document}

\title{Genetic Algorithm for Machine Scheduling\\
{\footnotesize 基於基因演化演算法解決機器排程問題}
% \thanks{Identify applicable funding agency here. If none, delete this.}
}


\author{\IEEEauthorblockN{1\textsuperscript{st} Hung Ming-Chun}
\IEEEauthorblockA{\textit{Department of Computer Science and} \\
\textit{Information Engineering}\\
\textit{National Taiwan Normal University}\\
Taipei, Taiwan(R.O.C) \\
61447026s@ntnu.edu.tw}
\and
\IEEEauthorblockN{2\textsuperscript{nd} Hsieh Yao-yu}
\IEEEauthorblockA{\textit{Graduate Institute of} \\
\textit{Science Education}\\
\textit{National Taiwan Normal University}\\
Taipei, Taiwan(R.O.C) \\
61345001s@ntnu.edu.tw}
\and
\IEEEauthorblockN{3\textsuperscript{rd} Given Name Surname}
\IEEEauthorblockA{\textit{Department of Computer Science and} \\
\textit{Information Engineering}\\
\textit{National Taiwan Normal University}\\
Taipei, Taiwan(R.O.C) \\
@ntnu.edu.tw}
}



\maketitle

\begin{abstract}

\end{abstract}

\section{Introduction}
在傳統直接式演算法中，
在處理NP問題時會遇到非常嚴重的複雜度瓶頸，
而啟發式演算法可以透過迭代，
能在固定時間下取得十分近似正確解答案，
而該研究利用基因演算法，透過不同的交配策略與突變機制，
並且在最後嘗試加入Tabu Search的記憶機制，來解決機器排程問題。



%%%%%%%%%%%%%%%%%%%%%%%%%%%%%%%%%%%%%%%%%%%%%%%%%%

\section{Related Work}
%這邊著重在介紹演算法，實作細節在Implement Detail
\subsection{Scheduling Problem}
% {\fontsize{10pt}{12pt}\selectfont
定序流線型工廠排程問題 (Permutation Flowshop Scheduling Problem, PFSP)
旨在決定 $n$ 個工作在 $m$ 台機器上的最佳加工順序。

本實驗之輸入包含
工作集合$J = \{j_1, j_2, \dots, j_n\}$、
機器集合 $M = \{M_1, M_2, \dots, M_m\}$
以及加工時間矩陣 $p_{i,j}$。

核心限制條件如下：
所有工作在各台機器上的加工順序 $\pi$ 必須完全一致，
且機器在同一時間僅能處理單一工作，
加工過程不可中斷。
目標函數為極小化總完工時間(Makespan)，記作 $C_{max}$。

定義 $C(\pi_k, i)$ 為位置 $k$ 的工作在機器 $i$ 的完工時間，
其遞迴關係如下：
\begin{equation}
    C(\pi_k, i) = \max\{C(\pi_{k-1}, i), C(\pi_k, i-1)\} + p_{i,\pi_k}
\end{equation}
最終目標為尋求一排列 $\pi^*$ 使得 $Z = C(\pi_n, m)$ 達最小值。
% }

本實驗之演算法實作包含以下兩個核心要素：
\begin{itemize}
    \setlength{\itemsep}{0pt}
    \item \textbf{解表示法 (Representation)}：採工作序列編碼。若有 $n$ 個工作，則解為一長度為 $n$ 的排列向量。
    \item \textbf{鄰域函式 (Neighborhood Function)}：採用交換移動 (Swap)。定義鄰域運算為隨機選取序列中兩個位置並交換其內容，以產生新的可行解。
\end{itemize}

\section{Genetic Algorithm}









\subsection{Tabu Search}
禁忌搜尋法利用記憶機制避免重複搜尋。
其\textbf{禁忌名單 (Tabu List)} 儲存的近代的排列方式，
旨在防止搜尋過程回流至已探索區域。
此外，本實作包含\textbf{特赦條件 (Aspiration Criterion)}：
若某一受禁忌限制的移動能產生優於目前歷史最佳解 $f_{best}$ 的結果，則無視禁忌狀態強制接受該移動。
% }


%%%%%%%%%%%%%%%%%%%%%%%%%%%%%%%%%%%%%%%%%%%%%%%%%%
\section{Implement Details}

\subsection{Scheduling Problem}
本研究針對排列流水線排程問題（Permutation Flowshop Scheduling Problem）設計了 \texttt{Scheduling} 類別作為評估函數（Evaluator）。給定一組工作的排列順序（Order），系統會依序將工作排入機器中。假設共有 $m$ 台機器，系統會維護一個大小為 $m$ 的陣列 \texttt{machine\_end\_time} 來記錄每台機器的最早可用時間。

對於排列中的每一個工作，其在第 $m_i$ 台機器的開始時間取決於「該工作在前一台機器的完工時間」與「第 $m_i$ 台機器的最早可用時間」兩者中的最大值：
$$start\_time =$$
$$ \max(current\_job\_ready\_time, machine\_end\_time[m_i])$$
接著將 $start\_time$ 加上該工作在該機器的處理時間 $job\_mtx[m_i][job]$ 即為完工時間。依此邏輯計算至最後一台機器，最後一個工作的完工時間即為該排程的總完工時間（Makespan）。

\subsection{Genetic Algorithm}

基因演算法（Genetic Algorithm, GA）是一種模擬自然界演化過程的群體導向元啟發式演算法。與單一軌跡搜尋不同，GA 維護一組族群（Population），透過個體間的資訊交換與競爭來探索複雜的解空間。

\begin{enumerate}
    \item \textbf{編碼與初始化：}
    在本實作中，染色體（Chromosome）採用排列編碼（Permutation Encoding），直接以工作的排列順序表示一個解。初始化階段利用 \texttt{generate\_random\_individual} 函數，透過 \texttt{std::shuffle} 產生 $N$ 個隨機排列的工作序列，並計算其初始完工時間（Makespan）作為適應度數值。

    \item \textbf{選汰機制（Selection）：}
    演算法採用錦標賽選擇法（Tournament Selection）。每次從族群中隨機抽選 $k$ 個個體（實作中預設 $k=3$），並從中選擇 Makespan 最小（適應度最高）的個體作為父母親（Parents）。此外，實作中加入了「菁英保留策略（Elitism）」，確保每一代中表現最優異的個體能直接晉級至下一代，防止優良基因因交配或突變而流失。

    \item \textbf{演化運算子（Genetic Operators）：}
    為了處理排列問題的合法性，系統實作了四種專用的交配運算子（Crossover）：順序交配（OX）、線性順序交配（LOX）、部分映射交配（PMX）與循環交配（CX）。突變算子則採用交換突變（Swap Mutation），隨機挑選染色體上的兩個位置進行交換，以維持族群的隨機搜尋能力。

    \item \textbf{禁忌強化與多樣性維持（Hybrid Tabu-GA）：}
    本實作最核心的特色在於整合了禁忌搜尋（Tabu Search）的防重複機制。當 \texttt{with\_tabu} 參數開啟時，演算法會執行以下進階操作：
    \begin{itemize}
        \item \textbf{狀態禁忌：} 
        實作中採用了多項式滾動雜湊（Polynomial Rolling Hash）來計算每個排列的 Hash 值：
        $$H = \sum_{k=0}^{n-1} (arr[k] + 1) \cdot p^k \pmod q$$
        其中 $arr[k]$ 為排列中第 $k$ 個位置的工作 ID，質數 $p$ 設為 $313$。模數 $q$ 則巧妙利用 C++ 中 \texttt{uint64\_t} 型態的自然溢位特性，等同於 $q = 2^{64}$。每次交配與突變產生新的排列狀態後，會檢查其 Hash 值是否在禁忌期間內。若是合法子代，則將其 Hash 值存入 \texttt{std::unordered\_map} 並記錄解禁代數（即當前代數加上禁忌長度）。


        \item \textbf{同代排除：} 在產生下一代的過程中使用 \texttt{std::unordered\_set} 比對 Hash，強制確保同一族群內不存在完全相同的雙胞胎個體。
        \item \textbf{隨機移民（Random Immigration）：} 當交配與突變經過 10 次嘗試（Retries）仍無法產生出不與禁忌名單衝突的獨特個體時，演算法會強制產生一個全新的隨機個體加入族群。這種機制能極大地延緩族群的「早熟收斂（Premature Convergence）」，在搜尋後期仍能保持優異的探索動能。
    \end{itemize}

    \item \textbf{監控指標：}
    為了量化演化過程，系統即時記錄了每一代的歷史最佳解（Global Best）、當代最佳解（Generation Best）以及族群完工時間的中位數（Median）。透過觀察中位數與最佳解的貼合程度，可以直觀地分析族群的收斂速率與基因同質化程度。


\end{enumerate}






\subsection{Order crossover (OX)}
順序交配（Order Crossover）旨在保證子代能形成合法的排列，同時保留基因的相對順序。
實作上，首先隨機選擇兩個切點（Cut points），並將第一個父代在切點範圍內的基因片段原封不動地複製給子代。
接著，從第二個父代的第二個切點後方開始遍歷，找出尚未被加入子代的剩餘基因，
並依序從子代的第二個切點後方開始填補。若遇到陣列尾端，則繞回（Wrap-around）陣列開頭繼續填入，直到填滿為止。

\subsection{Linear order crossover (LOX)}
線性順序交配（Linear Order Crossover）與 OX 類似，同樣先隨機選擇兩個切點並保留第一個父代的基因片段。
兩者最大的差異在於處理剩餘基因的「方向性」：LOX 捨棄了繞回機制，它會先收集第二個父代中尚未被使用的基因，
然後嚴格按照「由左至右」的順序，依序填入子代的空位中（直接跳過中間已複製的切點範圍），
從而更強烈地保留基因的絕對線性結構。

\subsection{Partially mapped crossover (PMX)}
部分映射交配（Partially Mapped Crossover）同時保留了基因的絕對位置與相對順序。
首先隨機選擇兩個切點，將父代的區段直接複製到對應的子代位置中，
並以此區段建立雙向的映射關係（Mapping dictionaries）。在填補切點以外的剩餘位置時，
子代會優先繼承另一個父代在該位置的基因；若該基因已存在於剛才複製的中間段落而發生衝突，
則利用建立好的映射表不斷追溯（實作中利用 \texttt{while} 迴圈解決連鎖映射問題），
直到找出一個未發生衝突的合法基因為止。

\subsection{Cycle crossover (CX)}
循環交配（Cycle Crossover）極端地追求「絕對位置」的保留。
其核心思想是確保子代每一個位置的基因，必定與其中一個父代在該位置的基因完全相同。
實作中透過追蹤兩個父代相同位置的基因形成「循環（Cycle）」。
位於循環內的基因直接繼承自第一個父代，而不在循環內的其餘基因則直接由第二個父代對應位置填補，
藉此保證產生合法的排列且不破壞絕對位置特徵。


\subsection{Tabu Mechanism in Genetic Algorithm}
本實作將禁忌機制（Tabu Mechanism）整合入基因演算法中，以維持族群多樣性並防止早熟收斂。在每一代產生新子代時，會利用禁忌名單過濾掉近期曾出現的排列解。為了高效比對狀態，實作中採用了「整體的排列狀態」作為禁忌對象：


\subsection{Gnuplot Visualizer}
為了直觀分析演算法的收斂行為與跳脫區域最佳解的過程，本實作內建了基於 \texttt{Gnuplot} 的視覺化工具。程式透過 \texttt{popen} 開啟 \texttt{gnuplot} 的行程通訊（Pipe），並利用標準輸入字元 \texttt{'-'} 將 C++ 陣列中的 \texttt{best\_makespan\_history} 與 \texttt{makespan\_history} 即時傳遞給繪圖引擎。

在大量平行測試的環境下，系統會透過 \texttt{save\_plot} 函數在背景靜默生成圖表。該函數會自動切換 Gnuplot 的 \texttt{terminal} 為 \texttt{jpeg}，並根據演算法名稱、測資名稱及最佳解數值自動命名檔案存入指定目錄，確保實驗數據的可重現性與易讀性。

%%%%%%%%%%%%%%%%%%%%%%%%%%%%%%%%%%%%%%%%%%%%%%%%%%%

\section{Experiment}
\fontsize{10pt}{12pt}\selectfont

本研究針對 Taillard 測試案例進行實驗，每個案例重複執行 20 次以確保統計穩定性。
\subsection{Different Crossover Operators}

\clearpage
\subsection{Different mutation rates for GA}
% 不同突變率的差別


\subsection{With or without Tabu Mechanism in GA}
\subsection{With or without Tabu Mechanism in GA}
% 比較有無 Tabu 機制的 GA (OX) 效能差異

為了驗證將禁忌機制（Tabu Mechanism）融入基因演算法的效益，我們以順序交配（OX）為基礎，分別測試了「純粹 GA」與「結合 Tabu 防重複機制之 GA」的表現。實驗數據如表 \ref{tab:ga_tabu_comp} 所示。

\begin{table}[htbp]
\caption{GA (OX) 與 GA (OX) 加上 Tabu 機制之結果比較}
\label{tab:ga_tabu_comp}
\centering
\resizebox{\columnwidth}{!}{
\begin{tabular}{l|ccc|ccc}
\hline
\textbf{Instance} & \multicolumn{3}{c|}{\textbf{GA (OX)}} & \multicolumn{3}{c}{\textbf{GA (OX) with Tabu}} \\
 & Min & Avg & Max & Min & Avg & Max \\ \hline
tai20\_5\_1   & 1297 & \textbf{1297.40} & 1305 & \textbf{1278} & 1298.15 & 1323 \\
tai20\_10\_1  & 1621 & \textbf{1641.10} & 1686 & \textbf{1605} & 1645.45 & 1677 \\
tai20\_20\_1  & \textbf{2321} & \textbf{2361.50} & 2400 & 2344 & 2390.85 & 2433 \\ \hline
tai50\_5\_1   & 2735 & 2748.60 & 2774 & \textbf{2724} & \textbf{2738.75} & 2752 \\
tai50\_10\_1  & 3190 & 3223.85 & 3288 & \textbf{3122} & \textbf{3171.75} & 3245 \\
tai50\_20\_1  & 4110 & 4182.50 & 4237 & \textbf{4026} & \textbf{4095.55} & 4212 \\ \hline
tai100\_5\_1  & 5505 & 5541.95 & 5564 & \textbf{5493} & \textbf{5508.55} & 5529 \\
tai100\_10\_1 & 6025 & 6108.55 & 6265 & \textbf{5907} & \textbf{5982.30} & 6050 \\
tai100\_20\_1 & 6770 & 6886.30 & 6978 & \textbf{6597} & \textbf{6706.75} & 6792 \\ \hline
\end{tabular}
}
\end{table}

\textbf{整體效益分析：大問題中的多樣性優勢} \\
觀察 \texttt{tai50} 與 \texttt{tai100} 系列的結果，可以明確看出 Tabu 機制帶來了壓倒性的優勢。無論是最佳解（Min）還是平均解（Avg），加入 Tabu 的 GA 皆顯著優於原始 GA。這印證了先前的理論：在解空間極度龐大的情況下，傳統 GA 容易發生近親繁殖與早熟收斂；而 Tabu 機制透過 Hash 比對強制剔除重複個體，成功迫使演算法持續探索未知的解空間，帶來了極佳的全域搜尋能力。

\textbf{反常現象探討：小問題中的過度強制多樣化} \\
然而，在小規模測資（\texttt{tai20} 系列）中，我們觀察到了截然不同的現象。特別是在 \texttt{tai20\_20\_1} 中，加入 Tabu 機制後的最差、平均乃至最佳解，反而全面劣於原始 GA。此外，在 \texttt{tai20\_5\_1} 與 \texttt{tai20\_10\_1} 中，Tabu 雖然幸運地探索到了更低的最佳解（Min），但其平均表現（Avg）卻變得較差且不穩定。

分析此反常現象的原因，在於\textbf{解空間規模與開發（Intensification）需求的權衡}。對於只有 20 個工作的排程問題，其優良解所在的「盆地（Basin of attraction）」相對狹窄且集中。原始 GA 能夠快速在該區域進行密集的交配與局部最佳化；反觀 Tabu 機制因為嚴格禁止生成曾出現的解，當族群逼近區域最佳解時，演算法會因為頻繁踩到「禁忌」而被迫進行「隨機移民（Random Immigration）」。這種\textbf{強制的過度多樣化（Over-diversification）}，反而硬生生地將演算法從快要收斂的優良區域中「踢走」，導致它浪費運算資源去探索劣質解，進而拖垮了整體的平均完工時間。

總結來說，Tabu 機制是防止大型問題早熟收斂的一劑猛藥，但對於解空間較小、需要高度密集開發（Intensification）的場景，適度關閉或降低 Tabu 的嚴格度，才能獲得更穩定的搜尋表現。


\clearpage

\section{Conclusion}
本研究成功實作並比較了 II、SA 與 TS 三種啟發式演算法於定序流線型工廠排程問題之表現。
透過大量的實驗數據分析，得出以下主要結論：
首先，在演算法對比中，TS 憑藉其獨特的記憶機制，在大型測資（如 tai100 系列）中顯著優於受限於貪婪策略的 II 以及需精細調參的 SA。
其次，在 SA 參數研究中，證實了較慢的降溫率雖增加計算代價，但能有效提升解的品質。
最重要的發現在於禁忌搜尋記憶內容的影響。實驗證明，相較於記錄完整解狀態的「排列狀態禁忌」，
「動作禁忌」具備更強的搜尋多樣化 (Diversification) 能力。由於排列狀態禁忌僅能精確避開曾出現過的單點解，
容易在相似結構的鄰域間產生震盪；而動作禁忌透過封鎖特定的交換路徑，能有效強制演算法探索全新的解空間。
未來的研究方向可考慮結合 SA 的接受機率與 TS 的禁忌名單，
開發混合型啟發式演算法 (Hybrid Metaheuristics)，
或引入動態禁忌長度調整機制，以進一步提升在極大規模工業排程問題上的運算效率。

% \begin{thebibliography}{00}
% \bibitem{b7} M. Young, The Technical Writer's Handbook. Mill Valley, CA: University Science, 1989.
% \end{thebibliography}

\end{document}
